\uuid{n6QF}
\exo7id{6902}
\titre{exo7 6902}
\auteur{ruette}
\organisation{exo7}
\datecreate{2013-01-24}
\isIndication{false}
\isCorrection{false}
\chapitre{Probabilité discrète}
\sousChapitre{Variable aléatoire discrète}
\cpge{mpsi-mp2i.probabilites-univers-fini.d, pcsi.probabilites-variables-aleatoires.d, ptsi.probabilites-variables-aleatoires.d, bcpst1.concepts-base-probabilites.a, bcpst1.variable-aleatoire-univers-fini.b, bcpst1.variable-aleatoire-univers-fini.d, mp-mpi.variables-aleatoires-discretes.c, pc.probabilites-lois-usuelles.c, psi.probabilites-lois-usuelles.c, pt.probabilites-lois-usuelles.d, bcpst2.variables-aleatoires-discretes.a}
\module{Probabilité et statistique}
\niveau{L2}
\difficulte{}

\contenu{
\question{
Soit $n$ un entier strictement positif.
Quelle est la loi du nombre de garçons dans une famille de $n$
enfants ? (Préciser les hypothèses que vous faites)
}
}
